%%% =======================================================================
%%% ISLS Extension Phase 3 v1.0.0
%%% C18: Hierarchical Multi-Scale Observation Layer
%%% Absorbing CubeTube Kybernetic Core Architecture
%%% Author: Sebastian Klemm
%%% Date: 16 March 2026
%%% Status: Implementation Specification for Claude Code
%%% =======================================================================
\documentclass[11pt,a4paper,twoside]{article}

\usepackage[utf8]{inputenc}
\usepackage{amsmath,amssymb,amsthm,mathtools}
\usepackage{algorithm,algorithmic}
\usepackage{booktabs,longtable,tabularx}
\usepackage{enumitem}
\usepackage{geometry}
\usepackage{hyperref}
\usepackage{cleveref}
\usepackage{xcolor}
\usepackage{fancyhdr}
\usepackage{tikz}
\usetikzlibrary{arrows.meta,positioning,shapes.geometric,fit,calc}

\geometry{margin=2.5cm,top=3cm,bottom=3cm,headheight=14pt}

\theoremstyle{definition}
\newtheorem{definition}{Definition}[section]
\newtheorem{axiom}{Axiom}[section]
\newtheorem{requirement}{Requirement}[section]
\newtheorem{invariant}{Invariant}[section]

\theoremstyle{plain}
\newtheorem{theorem}{Theorem}[section]
\newtheorem{proposition}{Proposition}[section]
\newtheorem{lemma}{Lemma}[section]
\newtheorem{corollary}{Corollary}[section]

\theoremstyle{remark}
\newtheorem{remark}{Remark}[section]
\newtheorem{example}{Example}[section]

\newcommand{\ISLS}{\textsf{ISLS}}
\newcommand{\RR}{\mathbb{R}}
\newcommand{\CC}{\mathbb{C}}
\newcommand{\NN}{\mathbb{N}}
\newcommand{\ZZ}{\mathbb{Z}}
\newcommand{\Hfive}{H^{5}}
\newcommand{\calL}{\mathcal{L}}
\newcommand{\calG}{\mathcal{G}}
\newcommand{\calR}{\mathcal{R}}
\newcommand{\calS}{\mathcal{S}}
\newcommand{\calK}{\mathcal{K}}
\newcommand{\calP}{\mathcal{P}}
\newcommand{\calT}{\mathcal{T}}
\newcommand{\calC}{\mathcal{C}}
\newcommand{\calH}{\mathcal{H}}
\newcommand{\calB}{\mathcal{B}}
\newcommand{\calM}{\mathcal{M}}
\newcommand{\calU}{\mathcal{U}}
\newcommand{\calD}{\mathcal{D}}
\DeclareMathOperator{\Tr}{Tr}
\DeclareMathOperator{\diag}{diag}
\DeclareMathOperator{\spec}{spec}
\DeclareMathOperator{\Dig}{Dig}
\DeclareMathOperator{\Canon}{Canon}
\DeclareMathOperator{\Lift}{Lift}
\DeclareMathOperator{\Proj}{Proj}
\DeclareMathOperator{\Agg}{Agg}
\newcommand{\norm}[1]{\left\lVert #1 \right\rVert}
\newcommand{\abs}[1]{\left\lvert #1 \right\rvert}
\newcommand{\innerprod}[2]{\langle #1, #2 \rangle}

\pagestyle{fancy}
\fancyhf{}
\fancyhead[LE,RO]{\ISLS{} Extension Phase 3 v1.0.0}
\fancyhead[RE,LO]{Sebastian Klemm}
\fancyfoot[C]{\thepage}

\hypersetup{
  colorlinks=true,
  linkcolor=blue!60!black,
  pdftitle={ISLS Extension Phase 3 v1.0.0 — Multi-Scale Layer},
  pdfauthor={Sebastian Klemm},
}

\begin{document}

\begin{titlepage}
\centering
\vspace*{2cm}
{\Huge\bfseries \ISLS{} Extension Phase~3\\[0.3cm]
v1.0.0}\\[1.5cm]
{\Large C18: Hierarchical Multi-Scale\\
Observation Layer}\\[1cm]
{\large Hypercube Universes, Dimensional Bridges,\\
and Scale Ladders for the\\
Intelligent Semantic Ledger Substrate}\\[1cm]
{\normalsize Absorbing CubeTube Kybernetic Core\\
Architecture Patterns}\\[2cm]
{\large Sebastian Klemm}\\[0.5cm]
{\large 16 March 2026}\\[2cm]

\begin{abstract}
\noindent
This document specifies \texttt{isls-scale} (C18), a hierarchical
multi-scale observation layer for \ISLS{}. The crate enables the system
to observe, analyze, and crystallize patterns simultaneously across
three scales---micro (individual entities), meso (synchronization
clusters), and macro (system-of-clusters)---using a unified formalism
of Hypercube Universes, Dimensional Bridges, and Scale Ladders.

Each scale is modeled as an independent state space (Hypercube
Universe) with its own transition dynamics and policy. Bridges couple
Hypercubes at the same scale (mesoscopic field coupling). Ladders
connect Hypercubes across scales via Lift and Projection operators
(microscopic--macroscopic linkage). Crystals produced at any scale
carry a hierarchical topological signature referencing structures at
all three scales.

The architecture absorbs the StateSpace, Policy, Transition, Dynamics,
HyperBounds, and Blueprint abstractions from CubeTube Kybernetic Core,
adapted to \ISLS{}'s deterministic, append-only, replay-verified
substrate.

\medskip\noindent\textbf{Status:} Implementation Specification.\\
\textbf{Parent:} ISLS v1.0.0 + Extensions Phase 1 + Phase 2.\\
\textbf{Absorbs:} CubeTube Kybernetic Core (architectural patterns).\\
\textbf{Depends on:} C16 (\texttt{isls-topology}) for spectral
clustering and Kuramoto synchronization.
\end{abstract}
\end{titlepage}

\tableofcontents
\newpage


%%% =====================================================================
\part{Foundations}\label{part:foundations}

\section{Purpose}\label{sec:purpose}

\ISLS{} currently observes a flat graph of entities at a single scale.
All entities are peers; there is no notion of containment, hierarchy,
or scale separation. This limits the system in three ways:

\begin{enumerate}[nosep]
  \item Patterns that manifest only at aggregate scales (e.g., sector
        rotations, regime transitions) are invisible to micro-level
        constraint extraction.
  \item The O($n^2$) edge complexity of S-Scale (2000 entities, 70s per
        tick) demonstrates that flat observation does not scale.
  \item Cross-domain observation (multiple independent systems in one
        \ISLS{} instance) has no structural separation.
\end{enumerate}

C18 resolves all three by introducing a hierarchical scale architecture
in which the persistent graph is partitioned into \emph{Hypercube
Universes} at three scales, connected by \emph{Bridges} (intra-scale
coupling) and \emph{Ladders} (inter-scale aggregation/disaggregation).


\section{Axioms}\label{sec:axioms}

\begin{axiom}[Scale Separation]\label{ax:scalesep}
Observable structure exists at multiple scales simultaneously. No
single scale is sufficient to capture all discoverable constraint
programs.
\end{axiom}

\begin{axiom}[Scale Invariance of Substrate]\label{ax:scaleinv}
The PSP substrate --- state space $\Sigma_0 = \RR^5$, canonicalization,
Run Descriptor, Evidence Chain, Crystal format, Gate cascade --- is
identical at every scale. Only the \emph{population} of the state space
(which entities, which edges, which observations) differs.
\end{axiom}

\begin{axiom}[Directed Scale Hierarchy]\label{ax:directed}
The scale hierarchy is strictly ordered: $\text{micro} \prec
\text{meso} \prec \text{macro}$. Information flows upward via
aggregation (Lift) and downward via disaggregation (Projection).
No scale may depend on a scale above it for its internal dynamics.
\end{axiom}

\begin{axiom}[Bridge Locality]\label{ax:bridgelocal}
Bridges couple Hypercubes at the same scale only. Cross-scale coupling
is handled exclusively by Ladders.
\end{axiom}

\begin{axiom}[Crystal Scale Provenance]\label{ax:crystalprov}
Every crystal carries provenance indicating the scale at which it was
discovered: micro, meso, or macro. A meso-crystal \textbf{MAY}
reference micro-crystals as sub-evidence. A macro-crystal \textbf{MAY}
reference meso-crystals.
\end{axiom}


\section{Relationship to Existing Crates}\label{sec:relationship}

\begin{tabularx}{\textwidth}{l l X}
\toprule
\textbf{Depends on} & \textbf{Provides to C18} & \textbf{Mechanism} \\
\midrule
\texttt{isls-types} & $\RR^5$ state, Hash256 & Foundation types \\
\texttt{isls-persist} & PersistentGraph, edges & Micro-scale substrate \\
\texttt{isls-topology} (C16) & Spectral clusters, Kuramoto $r_k$ &
  Meso-scale cluster identification \\
\texttt{isls-extract} & inverse\_weave & Per-scale constraint
  extraction \\
\texttt{isls-consensus} & Gate cascade & Per-scale crystallization \\
\texttt{isls-scheduler} (C15) & Spiral scheduler & Scale-adaptive
  tick resolution \\
\bottomrule
\end{tabularx}


%%% =====================================================================
\part{Scale Architecture}\label{part:scale}

\section{The Three Scales}\label{sec:threescales}

\begin{definition}[Micro Scale $\calS_\mu$]\label{def:micro}
The micro scale is the existing \ISLS{} persistent graph $G = (V, E)$
where $V$ is the set of observed entities and $E$ the set of
co-occurrence / correlation edges. Each vertex $v \in V$ carries a
$\RR^5$ embedding. This is the raw observation layer.
\end{definition}

\begin{definition}[Meso Scale $\calS_m$]\label{def:meso}
The meso scale is a graph $G_m = (V_m, E_m)$ where each vertex
$c \in V_m$ is a \emph{cluster} of micro-entities and each edge
$e \in E_m$ is an inter-cluster coupling. Clusters are identified by
spectral analysis (Fiedler vector bisection, C16) and/or Kuramoto
synchronization (C16).
\end{definition}

\begin{definition}[Macro Scale $\calS_M$]\label{def:macro}
The macro scale is a graph $G_M = (V_M, E_M)$ where each vertex
$u \in V_M$ is a \emph{domain} or \emph{regime} and each edge
$e \in E_M$ is an inter-domain coupling. Domains are identified by
second-order clustering: clustering the meso-clusters.
\end{definition}

\begin{remark}
The three scales are not fixed at compile time. They are emergent: the
system discovers how many clusters exist at the meso level and how many
domains at the macro level. If the data has no meaningful cluster
structure, $G_m$ collapses to a single node and the meso/macro layers
are trivial.
\end{remark}


\section{Hypercube Universe}\label{sec:hypercube}

\begin{definition}[Hypercube Universe]\label{def:universe}
A Hypercube Universe $\calU$ is a tuple
\[
  \calU = \langle \mathtt{id}, \mathtt{scale}, G_\calU, \Pi_\calU,
          T_\calU, \calP_\calU, \calB_\calU, \Sigma_\calU \rangle
\]
where:
\begin{itemize}[nosep]
  \item $\mathtt{id}$: content-addressed identifier,
  \item $\mathtt{scale} \in \{\mu, m, M\}$: scale level,
  \item $G_\calU = (V_\calU, E_\calU)$: the state space graph (directed,
        weighted),
  \item $\Pi_\calU$: the set of $\RR^5$ embeddings for all vertices,
  \item $T_\calU$: the set of admissible transitions (CubeTube
        \texttt{Transition} trait),
  \item $\calP_\calU$: the active policy (CubeTube \texttt{Policy}
        trait),
  \item $\calB_\calU$: the 5D bounding box (\texttt{HyperBounds}) of
        the universe,
  \item $\Sigma_\calU$: aggregate state vector of the universe (a single
        $\RR^5$ point summarizing the collective state).
\end{itemize}
\end{definition}

\begin{definition}[Aggregate State]\label{def:aggstate}
The aggregate state of a Hypercube Universe is computed by an
aggregation operator $\Agg$:
\[
  \Sigma_\calU = \Agg(\Pi_\calU) =
  \left(
    \bar{p},\; \bar{\rho},\; \bar{\omega},\; \bar{\chi},\; r_\calU
  \right)
\]
where $\bar{p}, \bar{\rho}, \bar{\omega}, \bar{\chi}$ are the
weighted means of the entity embeddings (weighted by edge degree) and
$r_\calU$ is the Kuramoto order parameter of the universe's entities.
\end{definition}

\begin{definition}[HyperBounds]\label{def:hyperbounds}
A HyperBounds object defines the axis-aligned bounding box in $\RR^5$:
\[
  \calB = [\min_p, \max_p] \times [\min_\rho, \max_\rho] \times
          [\min_\omega, \max_\omega] \times [\min_\chi, \max_\chi] \times
          [\min_\eta, \max_\eta]
\]
with operations:
\begin{itemize}[nosep]
  \item $\mathtt{contains}(\sigma)$: test whether $\sigma \in \calB$,
  \item $\mathtt{volume}()$: 5D hypervolume
        $\prod_{d=1}^{5}(\max_d - \min_d)$,
  \item $\mathtt{split}(d)$: bisect along dimension $d$, producing two
        child HyperBounds,
  \item $\mathtt{split\_all}()$: produce $2^5 = 32$ children (full
        subdivision),
  \item $\mathtt{merge}(\calB')$: smallest enclosing box.
\end{itemize}
\end{definition}


\section{Transitions and Dynamics}\label{sec:transitions}

\begin{definition}[Scale Transition]\label{def:transition}
A scale transition is a deterministic map
$T: \Sigma_\calU \to \Sigma_\calU$
that evolves the aggregate state of a universe. The canonical
transition vocabulary is:
\begin{enumerate}[nosep,label=(\roman*)]
  \item $\mathtt{Identity}$: $\Sigma \mapsto \Sigma$,
  \item $\mathtt{Translate}(\delta)$: $\Sigma \mapsto \Sigma + \delta$
        for offset $\delta \in \RR^5$,
  \item $\mathtt{Scale}(\alpha)$: $\Sigma \mapsto \alpha \cdot \Sigma$,
  \item $\mathtt{PhaseAdvance}(\dot\omega)$: increment the $\omega$
        coordinate by $\dot\omega \cdot \Delta t$,
  \item $\mathtt{LerpToward}(\Sigma^*, \gamma)$:
        $\Sigma \mapsto \Sigma + \gamma(\Sigma^* - \Sigma)$,
  \item $\mathtt{Decay}(\tau)$:
        $\Sigma \mapsto \Sigma \cdot e^{-\Delta t / \tau}$.
\end{enumerate}
All transitions \textbf{MUST} be deterministic under fixed
$(\Sigma, \mathrm{RD})$.
\end{definition}

\begin{definition}[Scale Dynamics]\label{def:dynamics}
Scale dynamics extend transitions with time and external input:
\[
  \Sigma_{k+1} = f(\Sigma_k, I_k, \Delta t, \mathrm{RD})
\]
where $I_k$ is an external input signal (e.g., observations flowing into
the universe from a Ladder below). The dynamics function $f$
\textbf{MUST} be a registered operator in the registry (C12).
\end{definition}

\begin{definition}[Scale Policy]\label{def:policy}
A scale policy $\calP$ governs which transition is applied at each
step within a universe:
\begin{enumerate}[nosep,label=(\roman*)]
  \item $\mathtt{Cyclic}$: round-robin through a fixed transition
        sequence,
  \item $\mathtt{Greedy}(\Sigma^*)$: select the transition that
        minimizes $\norm{\Sigma_{k+1} - \Sigma^*}$,
  \item $\mathtt{Balanced}$: prefer less-visited regions of the
        state space (exploration bias),
  \item $\mathtt{Reactive}$: select transition based on current metric
        thresholds (e.g., switch to Decay when $r_\calU < \theta_r$),
  \item $\mathtt{Custom}(f)$: user-defined deterministic function.
\end{enumerate}
Policies \textbf{MUST} be deterministic. $\mathtt{Random}$ from
CubeTube is excluded because \ISLS{} forbids \texttt{rand}.
\end{definition}


%%% =====================================================================
\part{Bridges and Ladders}\label{part:bridgeladder}

\section{Dimensional Bridges}\label{sec:bridges}

\begin{definition}[Bridge]\label{def:bridge}
A bridge $\calB_{ij}$ is a directed coupling between two Hypercube
Universes $\calU_i$ and $\calU_j$ at the same scale:
\[
  \calB_{ij} = \langle \calU_i, \calU_j, w_{ij}, \tau_{ij},
               \varphi_{ij}, \Pi_{ij} \rangle
\]
where:
\begin{itemize}[nosep]
  \item $w_{ij} \in [0, 1]$: coupling strength (normalized),
  \item $\tau_{ij} \in \NN_{\ge 0}$: propagation delay in ticks,
  \item $\varphi_{ij} \in [0, 2\pi)$: phase offset,
  \item $\Pi_{ij}$: the bridge policy governing when the coupling is
        active.
\end{itemize}
\end{definition}

\begin{definition}[Bridge Activation]\label{def:bridgeact}
A bridge is active at tick $k$ if and only if
\[
  \Pi_{ij}(\Sigma_{\calU_i,k}, \Sigma_{\calU_j,k}, \theta_{\mathrm{bridge}})
  = 1.
\]
The canonical activation predicate is:
\[
  \Pi_{ij} = \mathbf{1}\!\left[
    \norm{\Sigma_{\calU_i} - \Sigma_{\calU_j}} \le \theta_{\mathrm{prox}}
    \;\wedge\;
    r_{\calU_i} \ge \theta_{r,\min}
    \;\wedge\;
    r_{\calU_j} \ge \theta_{r,\min}
  \right]
\]
i.e., the bridge activates when both universes are sufficiently coherent
and their aggregate states are sufficiently proximate.
\end{definition}

\begin{definition}[Bridge Effect]\label{def:bridgeeff}
When bridge $\calB_{ij}$ is active, the coupling modifies the target
universe's dynamics:
\[
  \Sigma_{\calU_j,k+1} = f_j(\Sigma_{\calU_j,k}, I_k, \Delta t)
  + w_{ij} \cdot g(\Sigma_{\calU_i, k - \tau_{ij}}, \varphi_{ij})
\]
where $g$ is the bridge transfer function:
\[
  g(\Sigma, \varphi) = \Sigma \cdot \cos(\omega(\Sigma) + \varphi)
\]
extracting the phase-aligned component of the source universe's state.
\end{definition}

\begin{invariant}[Bridge Symmetry Breaking]\label{inv:bridgesym}
A bridge is directed: $\calB_{ij}$ does not imply $\calB_{ji}$.
Bidirectional coupling requires two explicit bridges. This preserves
causal directionality.
\end{invariant}


\section{Scale Ladders}\label{sec:ladders}

\begin{definition}[Ladder]\label{def:ladder}
A scale ladder $\calL_{s \to s'}$ is a pair of operators connecting
adjacent scales:
\[
  \calL_{s \to s'} = (\Lift_{s \to s'}, \Proj_{s' \to s})
\]
where $s \prec s'$ (e.g., $\mu \prec m$ or $m \prec M$).
\end{definition}

\subsection{Micro $\to$ Meso Ladder}

\begin{definition}[Micro-to-Meso Lift]\label{def:m2mlift}
Given a spectral cluster assignment $c: V \to V_m$ (from C16's Fiedler
bisection or Kuramoto synchronization), the lift operator constructs
the meso graph:
\[
  \Lift_{\mu \to m}(G, c) = G_m = (V_m, E_m, \Pi_m)
\]
where:
\begin{itemize}[nosep]
  \item $V_m = \{c_1, \ldots, c_K\}$ is the set of clusters,
  \item $\Pi_m(c_i) = \Agg(\{\Pi(v) : c(v) = c_i\})$ is the aggregate
        embedding of cluster $c_i$,
  \item $E_m(c_i, c_j) = \sum_{v \in c_i, w \in c_j} W_{vw}$ is the
        inter-cluster edge weight (sum of micro-edge weights between
        cluster members).
\end{itemize}
\end{definition}

\begin{definition}[Meso-to-Micro Projection]\label{def:m2mproj}
The projection operator disaggregates a meso-level signal $\delta_m$
back to the micro level:
\[
  \Proj_{m \to \mu}(\delta_m, c) =
  \left\{
    \frac{\delta_{c(v)}}{\abs{c(v)}} : v \in V
  \right\}
\]
distributing the meso-level perturbation uniformly to all members of
each cluster.
\end{definition}

\subsection{Meso $\to$ Macro Ladder}

\begin{definition}[Meso-to-Macro Lift]\label{def:meso2macro}
A second-order clustering $d: V_m \to V_M$ groups meso-clusters into
domains. The lift is structurally identical to $\Lift_{\mu \to m}$:
\[
  \Lift_{m \to M}(G_m, d) = G_M = (V_M, E_M, \Pi_M)
\]
with domain-aggregate embeddings and inter-domain edge weights.
\end{definition}

\begin{definition}[Macro-to-Meso Projection]\label{def:macro2meso}
Identical structure to $\Proj_{m \to \mu}$, distributing macro-level
signals to meso-clusters.
\end{definition}

\subsection{Ladder Consistency}

\begin{theorem}[Lift-Project Quasi-Inverse]\label{thm:liftproj}
$\Proj_{s' \to s} \circ \Lift_{s \to s'}$ is not the identity (it
loses intra-cluster variation) but satisfies:
\[
  \Agg(\Proj_{m \to \mu}(\Lift_{\mu \to m}(G, c))) = \Agg(G)
\]
i.e., the global aggregate state is preserved under the
lift-then-project round-trip.
\end{theorem}

\begin{proof}
$\Lift$ computes per-cluster weighted means. $\Proj$ distributes these
means uniformly. The global weighted mean of the projected values equals
the original global weighted mean because:
\[
  \frac{1}{\abs{V}} \sum_{v \in V} \frac{\Agg(c(v))}{\abs{c(v)}} \cdot \abs{c(v)}
  = \frac{1}{\abs{V}} \sum_{i} \Agg(c_i) \cdot \abs{c_i}
  = \Agg(G).
\qedhere
\]
\end{proof}

\begin{invariant}[Ladder Direction]\label{inv:ladderdir}
Information flows upward (Lift) passively and continuously. Information
flows downward (Projection) only when a crystal or alert at a higher
scale triggers a disaggregation signal. This prevents downward
information leakage from polluting micro-level observations.
\end{invariant}


%%% =====================================================================
\part{Multi-Scale Engine Integration}\label{part:engine}

\section{Per-Scale Macro-Step}\label{sec:perscale}

\begin{definition}[Scale-Partitioned Macro-Step]\label{def:scalemacro}
At each engine tick $k$, the multi-scale macro-step executes:
\begin{enumerate}[nosep,label=(MS\arabic*)]
  \item \textbf{Micro tick}: Execute the existing L0$\to$L4 pipeline
        on $G$, producing micro-embeddings, micro-constraints,
        micro-crystals.
  \item \textbf{Lift $\mu \to m$}: Construct $G_m$ from spectral
        clusters (C16) and micro-embeddings.
  \item \textbf{Bridge propagation (meso)}: For each active bridge
        $\calB_{ij}$ at meso scale, apply coupling effect.
  \item \textbf{Meso tick}: Run constraint extraction and gate cascade
        on $G_m$, producing meso-constraints, meso-crystals.
  \item \textbf{Lift $m \to M$}: Construct $G_M$ from meso-clusters.
  \item \textbf{Bridge propagation (macro)}: For each active bridge at
        macro scale, apply coupling effect.
  \item \textbf{Macro tick}: Run constraint extraction and gate cascade
        on $G_M$, producing macro-constraints, macro-crystals.
  \item \textbf{Downward projection (optional)}: If a macro-crystal
        triggers, project a disaggregation signal through
        $\Proj_{M \to m} \circ \Proj_{m \to \mu}$ to refine
        micro-level observation.
\end{enumerate}
\end{definition}

\begin{requirement}[Scale Budget]\label{req:scalebudget}
Each scale has an independent compute budget. If the meso or macro
scale exceeds its budget, it is skipped for that tick. The micro scale
is never skipped.
\end{requirement}

\begin{requirement}[Scale-Aware Scheduling]\label{req:scalesched}
The spiral scheduler (C15) \textbf{MAY} assign different sub-step
counts to different scales:
\begin{itemize}[nosep]
  \item Micro: $n_k^{(\mu)}$ sub-steps (existing behavior),
  \item Meso: $n_k^{(m)} = \max(1, n_k^{(\mu)} / \alpha_m)$ where
        $\alpha_m \ge 1$ is the meso decimation factor,
  \item Macro: $n_k^{(M)} = \max(1, n_k^{(\mu)} / \alpha_M)$ where
        $\alpha_M \ge \alpha_m$.
\end{itemize}
This ensures higher scales tick less frequently than lower scales.
\end{requirement}


\section{Cluster Identification}\label{sec:clustering}

\begin{definition}[Spectral Bisection Clustering]\label{def:speccluster}
Using the Fiedler vector $u_1$ from C16's spectral decomposition:
\begin{enumerate}[nosep]
  \item Sort vertices by their Fiedler vector component $u_1(v)$.
  \item Bisect at the median: $c_0 = \{v : u_1(v) \le \tilde{u}_1\}$,
        $c_1 = \{v : u_1(v) > \tilde{u}_1\}$.
  \item Recursively bisect each half until the spectral gap of a
        sub-cluster falls below $\theta_{\Delta,\min}$ or
        $\abs{c_i} < n_{\min}$.
\end{enumerate}
\end{definition}

\begin{definition}[Kuramoto Synchronization Clustering]\label{def:kuramcluster}
As a complementary method (or replacement when spectral decomposition
is budget-constrained):
\begin{enumerate}[nosep]
  \item Run Kuramoto integration for $N_{\mathrm{sync}}$ steps (C16).
  \item Group vertices whose phase difference
        $\abs{\phi_i - \phi_j} < \theta_\phi$ into the same cluster.
  \item Merge clusters smaller than $n_{\min}$ into their nearest
        neighbor cluster.
\end{enumerate}
\end{definition}

\begin{requirement}[Deterministic Clustering]\label{req:detcluster}
Both clustering methods \textbf{MUST} produce identical results under
identical $(\Sigma, \mathrm{RD})$. Tie-breaking \textbf{MUST} use the
canonical vertex ID ordering.
\end{requirement}


\section{Multi-Scale Crystallization}\label{sec:mscrystal}

\begin{definition}[Scale-Tagged Crystal]\label{def:scalecrystal}
A crystal produced by the multi-scale engine carries an extended
provenance:
\[
  C = \langle \ldots, \mathtt{scale}, \mathtt{universe\_id},
      \mathtt{sub\_crystal\_ids}, \mathtt{parent\_crystal\_ids},
      \tau_{\mathrm{multi}} \rangle
\]
where:
\begin{itemize}[nosep]
  \item $\mathtt{scale} \in \{\mu, m, M\}$: the scale at which this
        crystal was discovered,
  \item $\mathtt{universe\_id}$: the Hypercube Universe that produced
        it,
  \item $\mathtt{sub\_crystal\_ids}$: list of crystal IDs at the
        scale below that constitute evidence for this crystal,
  \item $\mathtt{parent\_crystal\_ids}$: list of crystal IDs at the
        scale above that reference this crystal,
  \item $\tau_{\mathrm{multi}}$: the multi-scale topological signature.
\end{itemize}
\end{definition}

\begin{definition}[Multi-Scale Topological Signature]\label{def:multitopo}
The multi-scale topological signature extends the hardened signature
from C16:
\[
  \tau_{\mathrm{multi}}(C) = \langle
    \tau_\mu, \tau_m, \tau_M,
    n_{\mathrm{clusters}}, n_{\mathrm{domains}},
    n_{\mathrm{bridges}}, n_{\mathrm{active\_bridges}},
    \mathrm{scale\_coherence}
  \rangle
\]
where:
\begin{itemize}[nosep]
  \item $\tau_\mu, \tau_m, \tau_M$: per-scale topological signatures
        (from C16),
  \item $n_{\mathrm{clusters}}$: number of meso-clusters,
  \item $n_{\mathrm{domains}}$: number of macro-domains,
  \item $n_{\mathrm{bridges}}$: total bridges,
  \item $n_{\mathrm{active\_bridges}}$: bridges currently active,
  \item $\mathrm{scale\_coherence} = \frac{r_M + r_m + r_\mu}{3}$:
        average Kuramoto coherence across scales.
\end{itemize}
\end{definition}

\begin{theorem}[Cross-Scale Crystal Hierarchy]\label{thm:crystalhier}
If a macro-crystal $C_M$ references meso-crystals
$\{C_{m,1}, \ldots, C_{m,k}\}$ and each $C_{m,i}$ references
micro-crystals $\{C_{\mu,i,1}, \ldots\}$, then the entire hierarchy
is independently verifiable: each crystal at each scale passes the
8-check formal validation independently, and the sub-crystal references
are content-addressed.
\end{theorem}

\begin{proof}
By construction: each crystal is produced by the same staged closure
and gate cascade, regardless of scale. Content-addressing ensures
reference integrity. The 8-check validation (content\_address,
dual\_consensus, evidence\_chain, free\_energy, gate\_kairos,
immutability, operator\_versions, por\_trace) operates on the crystal's
own fields, not on its scale. Therefore each level validates
independently.
\end{proof}


%%% =====================================================================
\part{New Metrics}\label{part:metrics}

\section{Multi-Scale Metrics M28--M32}\label{sec:newmetrics}

\begin{definition}[M28 Cluster Count]\label{def:m28}
$M_{28} = |V_m|$: number of meso-clusters. Stability alert if
$M_{28}$ changes by more than 50\% between consecutive ticks
(cluster structure unstable).
\end{definition}

\begin{definition}[M29 Bridge Activity]\label{def:m29}
$M_{29} = n_{\mathrm{active\_bridges}} / n_{\mathrm{bridges}}$:
fraction of active bridges. Low values indicate decoupled universes;
high values indicate strong inter-cluster coupling.
\end{definition}

\begin{definition}[M30 Scale Coherence]\label{def:m30}
$M_{30} = (r_M + r_m + r_\mu) / 3$: average Kuramoto coherence across
all three scales. Values near 1 indicate that structure is coherent at
all levels; values near 0 indicate scale-dependent incoherence.
\end{definition}

\begin{definition}[M31 Lift Compression Ratio]\label{def:m31}
$M_{31} = |V_m| / |V|$: the ratio of meso-vertices to micro-vertices.
This measures how much the hierarchy compresses. Typical healthy range:
$0.01 \le M_{31} \le 0.2$.
\end{definition}

\begin{definition}[M32 Cross-Scale Crystal Rate]\label{def:m32}
$M_{32}$: number of crystals with $\mathtt{sub\_crystal\_ids} \neq
\emptyset$ per unit time. This measures how often the system discovers
patterns that span multiple scales.
\end{definition}


%%% =====================================================================
\part{Configuration}\label{part:config}

\section{Scale Configuration}\label{sec:scaleconfig}

\begin{verbatim}
scale:
  enabled: true
  micro:
    # Micro scale uses existing engine config
  meso:
    enabled: true
    clustering_method: "spectral"   # or "kuramoto" or "hybrid"
    min_cluster_size: 5
    max_clusters: 50
    spectral_gap_threshold: 0.01
    phase_threshold: 0.3            # for Kuramoto clustering
    decimation_factor: 5            # alpha_m
    budget_ms: 50
  macro:
    enabled: true
    min_domain_size: 2
    max_domains: 10
    decimation_factor: 25           # alpha_M
    budget_ms: 20
  bridges:
    proximity_threshold: 0.5        # theta_prox
    coherence_threshold: 0.3        # theta_r_min
    max_delay_ticks: 10
  ladders:
    lift_method: "weighted_mean"    # or "median" or "kuramoto_order"
    projection_method: "uniform"    # or "weighted"
    downward_on_crystal: true       # project on macro crystal
\end{verbatim}


%%% =====================================================================
\part{Rust API}\label{part:api}

\section{Core Types}\label{sec:api:types}

\begin{verbatim}
// isls-scale/src/lib.rs

#[derive(Debug, Clone, Copy, PartialEq, Eq, Hash, Serialize, Deserialize)]
pub enum Scale { Micro, Meso, Macro }

#[derive(Debug, Clone, Serialize, Deserialize)]
pub struct HyperBounds {
    pub min: FiveDState,
    pub max: FiveDState,
}

impl HyperBounds {
    pub fn contains(&self, state: &FiveDState) -> bool;
    pub fn volume(&self) -> f64;
    pub fn split(&self, dimension: usize) -> (HyperBounds, HyperBounds);
    pub fn split_all(&self) -> [HyperBounds; 32];
    pub fn merge(&self, other: &HyperBounds) -> HyperBounds;
    pub fn from_points(points: &[FiveDState]) -> Option<HyperBounds>;
    pub fn center(&self) -> FiveDState;
}
\end{verbatim}

\section{Universe}\label{sec:api:universe}

\begin{verbatim}
#[derive(Debug, Clone, Serialize, Deserialize)]
pub struct HypercubeUniverse {
    pub id: Hash256,
    pub scale: Scale,
    pub vertex_ids: Vec<VertexId>,
    pub aggregate_state: FiveDState,
    pub bounds: HyperBounds,
    pub policy: ScalePolicy,
    pub kuramoto_r: f64,
}

#[derive(Debug, Clone, Serialize, Deserialize)]
pub enum ScalePolicy {
    Cyclic { sequence: Vec<String> },
    Greedy { target: FiveDState },
    Balanced,
    Reactive { thresholds: BTreeMap<String, f64> },
    Custom { name: String },
}

pub fn build_universe(
    vertices: &[VertexId],
    embeddings: &BTreeMap<VertexId, FiveDState>,
    scale: Scale,
    policy: ScalePolicy,
) -> HypercubeUniverse;

pub fn compute_aggregate(
    embeddings: &BTreeMap<VertexId, FiveDState>,
    degrees: &BTreeMap<VertexId, f64>,
    phases: &[f64],
) -> FiveDState;
\end{verbatim}

\section{Bridges}\label{sec:api:bridges}

\begin{verbatim}
#[derive(Debug, Clone, Serialize, Deserialize)]
pub struct Bridge {
    pub source_id: Hash256,       // source universe
    pub target_id: Hash256,       // target universe
    pub weight: f64,              // coupling strength [0,1]
    pub delay_ticks: u64,         // propagation delay
    pub phase_offset: f64,        // phase alignment
    pub active: bool,             // current activation state
}

pub fn evaluate_bridges(
    universes: &BTreeMap<Hash256, HypercubeUniverse>,
    bridges: &mut [Bridge],
    config: &ScaleConfig,
);

pub fn apply_bridge_coupling(
    target: &mut HypercubeUniverse,
    source_history: &[FiveDState],    // ring buffer of past states
    bridge: &Bridge,
    tick: u64,
);
\end{verbatim}

\section{Ladders}\label{sec:api:ladders}

\begin{verbatim}
pub fn lift_micro_to_meso(
    graph: &PersistentGraph,
    clusters: &BTreeMap<VertexId, usize>,    // vertex -> cluster_id
    embeddings: &BTreeMap<VertexId, FiveDState>,
) -> (Vec<HypercubeUniverse>, Vec<Bridge>);

pub fn lift_meso_to_macro(
    meso_universes: &[HypercubeUniverse],
    domains: &BTreeMap<usize, usize>,        // cluster_id -> domain_id
) -> (Vec<HypercubeUniverse>, Vec<Bridge>);

pub fn project_macro_to_meso(
    signal: &FiveDState,
    domain: &HypercubeUniverse,
    meso_universes: &[HypercubeUniverse],
    domain_assignment: &BTreeMap<usize, usize>,
) -> BTreeMap<Hash256, FiveDState>;

pub fn project_meso_to_micro(
    signals: &BTreeMap<Hash256, FiveDState>,
    clusters: &BTreeMap<VertexId, usize>,
    meso_universes: &[HypercubeUniverse],
) -> BTreeMap<VertexId, FiveDState>;
\end{verbatim}

\section{Clustering}\label{sec:api:clustering}

\begin{verbatim}
pub fn spectral_bisection_cluster(
    fiedler_vector: &[f64],
    vertex_ids: &[VertexId],
    spectral_gap_threshold: f64,
    min_cluster_size: usize,
    max_clusters: usize,
) -> BTreeMap<VertexId, usize>;

pub fn kuramoto_phase_cluster(
    phases: &[f64],
    vertex_ids: &[VertexId],
    phase_threshold: f64,
    min_cluster_size: usize,
) -> BTreeMap<VertexId, usize>;

pub fn hybrid_cluster(
    fiedler: &[f64],
    phases: &[f64],
    vertex_ids: &[VertexId],
    config: &ScaleConfig,
) -> BTreeMap<VertexId, usize>;
\end{verbatim}

\section{Multi-Scale Engine Hook}\label{sec:api:engine}

\begin{verbatim}
pub struct MultiScaleState {
    pub meso_universes: Vec<HypercubeUniverse>,
    pub macro_universes: Vec<HypercubeUniverse>,
    pub meso_bridges: Vec<Bridge>,
    pub macro_bridges: Vec<Bridge>,
    pub cluster_assignment: BTreeMap<VertexId, usize>,
    pub domain_assignment: BTreeMap<usize, usize>,
    pub meso_crystals: Vec<SemanticCrystal>,
    pub macro_crystals: Vec<SemanticCrystal>,
}

pub fn multi_scale_tick(
    micro_state: &GlobalState,
    scale_state: &mut MultiScaleState,
    spectral: &SpectralDecomposition,
    kuramoto: &KuramotoState,
    config: &ScaleConfig,
) -> MultiScaleTickResult;

pub struct MultiScaleTickResult {
    pub meso_crystals: Vec<SemanticCrystal>,
    pub macro_crystals: Vec<SemanticCrystal>,
    pub metrics: MultiScaleMetrics,
    pub events: Vec<ScaleEvent>,
}

pub struct MultiScaleMetrics {
    pub m28_cluster_count: u64,
    pub m29_bridge_activity: f64,
    pub m30_scale_coherence: f64,
    pub m31_lift_compression: f64,
    pub m32_cross_scale_crystal_rate: f64,
}

pub enum ScaleEvent {
    ClusterSplit { cluster_id: usize, into: Vec<usize> },
    ClusterMerge { clusters: Vec<usize>, into: usize },
    BridgeActivated { bridge_idx: usize },
    BridgeDeactivated { bridge_idx: usize },
    DownwardProjection { from_scale: Scale, crystal_id: Hash256 },
    FixpointReached { scale: Scale },
}
\end{verbatim}


%%% =====================================================================
\part{Acceptance Tests}\label{part:tests}

\section{Test Catalogue}\label{sec:tests}

\begin{longtable}{l l p{7cm}}
\toprule
\textbf{ID} & \textbf{Name} & \textbf{Procedure} \\
\midrule
\endfirsthead
\midrule
\endhead
\bottomrule
\endfoot
AT-SC1 & HyperBounds containment & Create bounds; test points inside
  and outside; verify correctness. \\
AT-SC2 & HyperBounds subdivision & Split\_all(); verify 32 children;
  verify no overlap; verify union = parent. \\
AT-SC3 & Aggregate state conservation & Compute Agg; lift; project;
  compute Agg again; verify equal. \\
AT-SC4 & Spectral clustering determinism & Cluster twice with same
  input; verify identical assignments. \\
AT-SC5 & Kuramoto clustering & Run Kuramoto to convergence; cluster
  by phase; verify synchronized entities are co-clustered. \\
AT-SC6 & Bridge activation & Set two universes with proximate states
  and high coherence; verify bridge activates. Set distant states;
  verify bridge inactive. \\
AT-SC7 & Bridge coupling effect & Activate bridge; verify target
  universe state is modified by source. \\
AT-SC8 & Ladder lift-project round-trip & Lift micro$\to$meso;
  project meso$\to$micro; verify global aggregate preserved. \\
AT-SC9 & Multi-scale crystal provenance & Run S-Basic with scale
  enabled; verify meso-crystals contain sub\_crystal\_ids referencing
  micro-crystals. \\
AT-SC10 & Scale-disabled passthrough & Set scale.enabled=false;
  verify behavior identical to pre-extension. \\
AT-SC11 & Budget fallback & Set meso budget to 1ms on large graph;
  verify meso tick is skipped and logged. \\
AT-SC12 & Cross-scale crystal validation & Validate a macro-crystal
  and all referenced sub-crystals; verify all pass 8-check. \\
AT-SC13 & Multi-scale metrics & Run 100 ticks with scale; verify
  M28--M32 are populated and non-zero. \\
AT-SC14 & Downward projection & Trigger a macro-crystal; verify
  downward projection event and micro-level signals. \\
AT-SC15 & Policy determinism & Run Greedy and Reactive policies;
  verify deterministic under same RD. \\
\end{longtable}


%%% =====================================================================
\part{Dependency Graph and Build Order}\label{part:deps}

\section{Updated Dependency Graph}

\begin{verbatim}
isls-types        <-- (all crates)
    ^
isls-registry     <-- (all operator crates)
    ^
isls-observe  <-- isls-persist <-- isls-topology (C16)
                       ^                ^
                       |          isls-extract
                       |                ^
                       |          isls-scale (C18: uses topology + persist)
                       |                ^
               isls-consensus <-- isls-carrier
                       ^
               isls-archive <-- isls-morph
                       ^
               isls-manifest
                       ^
               isls-capsule
                       ^
               isls-scheduler (uses scale events for adaptive ticking)
                       ^
               isls-engine (+ multi_scale_tick integration)
                       ^
               isls-store (C17)
                       ^
               isls-harness (+ scale tests)
                       ^
               isls-cli (+ scale metrics in reports)
\end{verbatim}

\section{Build Order}

\begin{enumerate}[nosep]
  \item \texttt{isls-scale} (C18): depends on \texttt{isls-types},
        \texttt{isls-persist}, \texttt{isls-topology} (C16).
        No new external dependencies.
  \item Engine integration: wire \texttt{multi\_scale\_tick} into the
        engine loop after the micro tick.
  \item Scheduler integration: feed scale events into scheduling
        function.
  \item Crystal format extension: add scale provenance fields.
  \item CLI/report integration: add M28--M32 to metric dashboard and
        HTML report.
  \item Harness: add AT-SC1 through AT-SC15.
\end{enumerate}

\section{Test Count}

\begin{tabular}{l r r}
\toprule
\textbf{Phase} & \textbf{New Tests} & \textbf{Cumulative} \\
\midrule
Core (C1--C9) & 130 & 130 \\
Phase 1 (C10--C15) & 31 & 161 \\
Phase 2 (C16--C17) & 20 & 181 \\
Phase 3 (C18) & 15 & $\ge$ 196 \\
\bottomrule
\end{tabular}


%%% =====================================================================
\part{Implementation Handoff}\label{part:handoff}

\section{Verification Sequence}

\begin{verbatim}
# Build
cargo build --workspace --release 2>&1 | grep -c warning  -> 0

# Test
cargo test --workspace -- --test-threads=1  -> >= 196 passed

# S-Basic with multi-scale
rm -rf ~/.isls
cargo run --bin isls --release -- ingest --adapter synthetic --scenario S-Basic
cargo run --bin isls --release -- run --mode shadow --ticks 100

# Verify multi-scale output
# - Meso universes created (cluster count > 1)
# - At least 1 meso-crystal with sub_crystal_ids
# - M28-M32 populated in metrics
# - Bridges evaluated (some active, some inactive)

cargo run --bin isls --release -- validate formal
# All crystals pass (micro + meso + macro)
\end{verbatim}

\section{Completion Criteria}

Phase 3 is complete when:
\begin{enumerate}[nosep]
  \item \texttt{isls-scale} compiles with zero warnings.
  \item All AT-SC tests pass.
  \item S-Basic produces crystals at micro and meso scales.
  \item M28--M32 metrics appear in reports.
  \item Scale disabled produces identical results to pre-extension.
  \item Multi-scale crystals reference sub-crystals and pass validation.
  \item Total tests $\ge 196$.
\end{enumerate}


\vfill
\begin{center}
\rule{0.5\textwidth}{0.4pt}\\[0.5cm]
{\small Generated for \ISLS{} Extension Phase 3 v1.0.0.\\
Hypercube Universes. Dimensional Bridges. Scale Ladders.\\
Deterministic, append-only, replay-verified.}
\end{center}

\end{document}
